Particles produce electromagnetic showers. When these showers are incident on the ECAL, they will spread laterally over several crystals. This tendency to spread energy across several crystals is increased by the scattering of particles off of the physical structure of the detectors in front of the ECAL and by the proclivity of the CMS experiment's strong magnetic field to spread energy along the $\phi$ polar direction as defined where the z-axis is in the anticlockwise beam direction. Clustering algorithms are therefore used to sum energy deposits in adjacent crystals belonging to the same electromagnetic shower. The clustering algorithm first identifies "basic clusters" corresponding to local maxima of the energy deposits. Basic clusters are then merged together to form a "supercluster", which is extended in $\phi$ to recover the radiated energy of the shower. Because of the differences between the geometric arrangement of the crystals in the barrel and endcap regions, a different clustering algorithm is used in each region \cite{CMS:2013lxn}. Clusters identified in the ECAL, and the energy and incident time of these clusters, are used in the particle-flow (PF) algorithm to reconstruct and identify each individual particle in an event.

The ECAL is primarily responsible for the reconstruction of electrons and photons. Initially, both photons and electrons are reconstructed as a single type of particle, that is called egamma (EG) candidate. Electrons candidates are then selected from the EG candidates by associating a track that was reconstructed in the silicon detector with an EG candidate. Photon candidates are selected from EG candidates that meet selection criteria derived from the transverse shower width, the hadronic to electromagnetic energy ratio from the hadronic calorimeter (HCAL) and ECAL, and cluster isolation. Photon candidates that are also selected as an electron candidate are vetoed from consideration \cite{CMS:2015myp}.

% the circumference of the detector perpendicularly to the beam direction, ($\phi$).